\section{Discussion}
\label{sec:discussion}

The results in Section~\ref{sec:evaluation} suggest that a cheap, non-learned
signal can meaningfully complement a learned ranking score rather than merely
duplicate it. Prompt length costs nothing beyond tokenization, is always
available regardless of how confident or well-calibrated the LTR predictor
is, and, unlike the predictor's score, cannot itself be affected by
distribution shift in the prompt distribution. This likely explains why the
composite scheduler's advantage over the LTR-only scheduler is more
pronounced on the out-of-distribution dataset than on the in-distribution
one: the length signal has the most to contribute exactly when the learned
score is least reliable.

\subsection{Limitations}
\label{subsec:limitations}

Several limitations qualify these findings. First, prompts in this evaluation
were submitted to the serving model without applying its chat template; we
verified in earlier diagnostic runs that this caused a substantial fraction of
generations to reach the configured maximum token cap rather than terminate
naturally, which compresses the observed spread of true output lengths and
should be considered when interpreting the absolute latency and HOL-blocking
values reported here. Second, the composite weights $\alpha$ and $\beta$ in
Eq.~\eqref{eq:composite-score} were fixed manually rather than learned or
tuned per dataset, so the reported results reflect one operating point rather
than the best achievable trade-off. Third, all scheduling policies in this
evaluation compute an admission order over the entire batch of requests
before submission, rather than reordering a continuously arriving stream in
real time; this differs from FCFS, which processes requests as they truly
arrive, and is a potential confound when comparing tail-latency behavior
across policies. Fourth, the evaluation is limited to a single model
(Llama-3.1-8B-Instruct), a single GPU (one NVIDIA A100), and two
1{,}000-request held-out sets, which may not capture the full range of
production traffic patterns or hardware configurations. Finally, the serving
environment uses a newer vLLM release (0.6.1.post2) than the original LTR
framework's target version (0.4.1), a substitution required because the
original version has no prebuilt wheel for the Python 3.12 environment used
here and because the platform's pre-installed PyTorch build is not
ABI-compatible with vLLM's compiled CUDA extensions; this limits direct
numerical comparison to results reported under the original framework's
software stack.

\subsection{Future Work}
\label{subsec:future-work}

Over the next six months, we intend to: (1) rerun the full benchmark suite
with the model's chat template correctly applied, to obtain output-length
distributions and latency figures free of the max-token-cap artifact
described above; (2) replace the fixed $\alpha, \beta$ weighting with a
learned or validation-tuned combination, potentially conditioned on the
predictor's own confidence; (3) evaluate the composite scheduler under a
true streaming arrival process rather than a pre-sorted batch, to remove the
admission-order confound noted above; (4) extend the evaluation to
additional model sizes and a multi-GPU, tensor-parallel serving
configuration; and (5) investigate using time-to-first-token directly as a
scheduling objective, rather than only as an evaluation metric, given its
sensitivity to queuing delay observed in Section~\ref{subsec:results-ttft}.
